% from ega4-17-fr.tex
\oldpage[IV-4]{78}

$\operatorname{Supp}(\operatorname{Ker}(1\otimes u_s))=(\operatorname{Supp}(\operatorname{Ker}(1\otimes u)))_s$ (I, 9.1.13.1); il résulte enfin de (9.5.2) que pour $s$ dans un voisinage de $\eta$, on a $(\operatorname{Supp}(\operatorname{Ker}(1\otimes u)))_s\supset T_s$, ce qui établit le lemme.

$3^\circ$ Supposons maintenant que $x\in X-E$, que le morphisme $f_\eta$ soit plat au point $x$, mais que $f_\eta$ ne soit pas lisse au point $x$. Remarquons que dire que $f_\eta$ est plat au point $x$ équivaut à dire que $f$ lui-même est plat au point $x$ et en remplaçant $X$ par un voisinage de $x$, on peut supposer que $f$ est plat (11.1.1); on en conclut qu'il en est de même de $f_s$ pour tout $s\in S$, et comme pour tout $y\in Y$, $f^{-1}(y)=f_{h(y)}^{-1}(y)$, il revient au même de dire que $f_{g(x')}$ est lisse au point $x'$ ou de dire que $f$ est lisse au point $x'$. Mais l'ensemble des $x'\in X$ où $f$ est lisse est \emph{ouvert} dans $X$ (12.1.7), donc l'ensemble des $x'\in X$ où $f$ n'est pas lisse est fermé, et puisqu'il contient $x$ par hypothèse, il contient aussi $\overline{\{x\}}$, ce qui achève la démonstration pour la première propriété considérée dans (i).

Prouvons en second lieu (i) lorsqu'il s'agit de la propriété d'être \emph{étale}. Notons pour cela que cette propriété pour $f_s$ au point $x$ équivaut à dire que $f_s$ est à la fois lisse et quasi-fini au point $x$ (17.6.1). Or, il revient au même de dire que $f_{g(x)}$ est quasi-fini au point $x$ ou que $f$ est lui-même quasi-fini en ce point; il résulte donc de (13.1.4) que l'ensemble des points $x$ tels que $f_{g(x)}$ soit quasi-fini au point $x$ est ouvert dans $X$, et \emph{a fortiori} localement constructible; la conclusion résulte donc de ce que l'ensemble des $x$ tels que $f_{g(x)}$ soit lisse au point $x$ est lui aussi localement constructible.

Passons à la preuve de (i) lorsqu'il s'agit de la propriété d'être \emph{différentiellement lisse}. Soit $p:X\times_Y X\to X$ la seconde projection canonique; la seconde projection $X_s\times_{Y_s}X_s\to X_s$ n'est autre que $p_s$ pour tout $s\in S$, et il résulte donc de (17.12.5.1)\footnote{Le lecteur vérifiera que le résultat de (17.7.11) n'est pas utilisé dans la suite du § 17 et qu'il n'y a donc pas de cercle vicieux.} que pour que $f_{g(x)}$ soit différentiellement lisse au point $x$, il faut et il suffit que $p_{g(x)}$ soit lisse au point $\Delta_f(x)$. Comme $p$ est localement de présentation finie, l'ensemble des points $z\in X\times_Y X$ tels que $p_{g(p(z))}$ soit lisse au point $z$ est localement constructible, et l'intersection de cet ensemble avec l'ensemble localement fermé $\Delta_f(X)$ est donc aussi localement constructible dans $\Delta_f(X)$ (1.8.2); la restriction de $p$ à $\Delta_f(X)$ étant un isomorphisme sur $X$, on voit que l'ensemble des $x\in X$ tels que $f_{g(x)}$ soit différentiellement lisse au point $x$ est localement constructible.

Considérons enfin la propriété d'être non ramifié; notons que le morphisme diagonal $\Delta_f:X\to X\times_Y X$ est une immersion localement de présentation finie (1.4.3.1) et pour tout $s\in S$, le morphisme diagonal $\Delta_{f_s}:X_s\to X_s\times_{Y_s}X_s$ n'est autre que $(\Delta_f)_s$; dire que $f_{g(x)}$ est non ramifié au point $x$ équivaut à dire que $(\Delta_f)_{g(x)}$ est un isomorphisme local au point $x$ (17.4.1), et puisque c'est une immersion localement de présentation finie, il revient au même (17.9.1)\footnote{Le lecteur vérifiera que le résultat de (17.7.11) n'est pas utilisé dans le reste du § 17 et qu'il n'y a donc pas de cercle vicieux.} de dire que $(\Delta_f)_{g(x)}$ est étale au point $x$; il suffit donc d'appliquer ce qui a été vu plus haut pour la propriété d'être étale.

